\section{Chapter~\ref{sec:dofaalt}. Other theories of \nt{DOPA}}

Each of the extended pictures rests on its own direct measurements of dopamine (the spread of reversal points, responses to unpleasant and novel events, the slow ramp, responses to a flavor switch), but the formulas that explain them are theories, and between two of them the experiments so far contradict one another.

{\footnotesize
\setlength{\tabcolsep}{3pt}\renewcommand{\arraystretch}{1.15}
\begin{longtable}{>{\raggedright}p{0.5cm}>{\raggedright}p{4.0cm}>{\raggedright}p{5.6cm}>{\raggedright}p{2.3cm}>{\raggedright\arraybackslash}p{1.7cm}}
\toprule
No. & Claim & Experiment and what was measured & Source & Status \\
\midrule\endfirsthead
\toprule
No. & Claim & Experiment and what was measured & Source & Status \\
\midrule\endhead
1 & The asymmetric rule~\eqref{eq:asymmetric} converges to the point~\eqref{eq:expectile}; for $\kappa=1/2$ it is the mean, and for a drop with probability $1/2$, $V=\kappa$ (Fig.~\ref{fig:expectiles}) & The point~\eqref{eq:expectile} is an expectile, the solution of a least-squares problem with different weights for positive and negative deviations; for the chapter's example the derivation is in the chapter itself. & \cite{NeweyPowell1987}; derived in the chapter & theory \\
2 & Different \nt{DOPA} neurons have different reversal points, ordered by asymmetry (Proposition~\ref{prop:expectile}) & Mice, optogenetically tagged \nt{DOPA} neurons of the ventral tegmental area, rewards of different sizes: the point where a burst turns into a dip differs between neurons; it is higher in neurons with greater response asymmetry; the shape of the reward distribution can be reconstructed from the group's responses. That this is the main function of the diversity is a hypothesis. & \cite{Dabney2020} & measured \\
3 & Some \nt{DOPA} neurons respond to unpleasant events with a burst & Monkeys: some \nt{DOPA} neurons are excited by a reward cue and inhibited by a cue for an air puff to the eye, others are excited by both; the groups lie in different parts of the midbrain. & \cite{Matsumoto2009, BrombergMartin2010} & measured \\
4 & The magnitude of the error, or novelty, sets the learning rate & The cited papers contain no direct experiment in which a salience signal in \nt{DOPA} neurons changes the learning rate; the review proposes the role of an ``attention, this matters'' signal. & \cite{BrombergMartin2010} & hypothesis \\
5 & A separate wire for the danger axis & Mice: \nt{DOPA} neurons projecting to the tail of the striatum respond to novel and intense stimuli, not to reward; ablating them weakens avoidance of a novel object. & \cite{Menegas2018} & measured \\
6 & The optimal action time $\tau_a^*=\sqrt{C/\bar R}$~\eqref{eq:vigor} & Minimum of the sum $C/\tau_a+\bar R\tau_a$; derived in the chapter. In the original model, the speed of actions is set by the cost of time, equal to the average reward. & \cite{Niv2007}; derived in the chapter & theory \\
7 & The slow dopamine level carries $\bar R$ and sets the speed of actions (Proposition~\ref{prop:multiplex}) & Rats, microdialysis and voltammetry in the nucleus accumbens: the dopamine level tracks the reward rate and the speed of actions from minute to minute. In humans, L-DOPA strengthens the dependence of reaction time on the average reward rate. That the slow level equals precisely $\bar R$ has not been proven. & \cite{Hamid2016, Beierholm2013vigor} & hypothesis \\
8 & The slow level combines two sources: release changes at the outputs without a change in spike rate, but a slow ramp also occurs in the rate itself & Rats, one task: release in the nucleus accumbens follows the changing reward expectation, while the spike rate of identified \nt{DOPA} neurons in the ventral tegmental area does not. Mice in a virtual corridor (row~10): a smooth ramp on approaching reward is present in the spikes of identified \nt{DOPA} neurons as well. & \cite{Mohebi2019, Kim2020} & measured \\
9 & Dopamine rises smoothly while the animal approaches a distant reward & Rats in a maze, voltammetry in the striatum: the dopamine concentration builds up over seconds as the goal approaches; the slope depends on the distance and on the reward size. & \cite{Howe2013ramp, Hamid2016} & measured \\
10 & The ramp is the error $\delta$ as the estimate is refined~\eqref{eq:ramp}; a jump in the corridor gives a burst (Fig.~\ref{fig:ramp}) & The formula and the value $\delta=0.75\,V$ per second are derived in the chapter. Experiment: mice in a virtual corridor, instantaneous teleportation closer to the goal; somatic spikes, somatic and axonal calcium, and the concentration in the striatum respond like an error, not like the expectation $V$. Which of the two explanations holds at all outputs is debated. & \cite{Kim2020} & measured \\
11 & Looking back: dopamine reports a measure of causal link~\eqref{eq:retro} & Mice, dopamine release in the nucleus accumbens in experiments where this theory and the error~\eqref{eq:ctd} predict different things: in several experiments release followed the former. & \cite{Jeong2022} & hypothesis \\
12 & The dopamine response gradually shifts from reward to cue during learning & Mice, signal from \nt{DOPA} neuron outputs in the ventral striatum over the course of learning: the burst gradually moves from the moment of reward to the moment of the cue, as learning by~\eqref{eq:ctd} requires. & \cite{Amo2022} & measured \\
13 & \nt{DOPA} neurons respond to a change in the reward's flavor at the same value & Rats, recordings from the ventral tegmental area: replacing a juice with another of equal value evokes a response, even though the error~\eqref{eq:ctd} is zero. & \cite{Takahashi2017} & measured \\
14 & Artificial bursts teach an association between two neutral stimuli & Rats, sensory preconditioning of two stimuli without reward: optogenetic activation of \nt{DOPA} neurons at the transition from the first stimulus to the second creates the association, and inhibition prevents it from being learned. & \cite{Sharpe2017} & measured \\
15 & A vector of errors per feature~\eqref{eq:vectortd} & A theory of feature-prediction errors; there is no direct test of the vector form. & \cite{Gardner2018} & hypothesis \\
16 & Different \nt{DOPA} neurons in the same task carry different quantities & Mice in a virtual corridor, two-photon imaging of \nt{DOPA} neurons: some are related to reward, others to speed and turning, still others to cues and choice accuracy. & \cite{Engelhard2019} & measured \\
17 & A bundle instead of a wire (Proposition~\ref{prop:bundle}) & Summary of rows~2--16: each decomposition is measured separately; there is no experimental evidence for a single picture in which all of them are parts of one signal. & --- & hypothesis \\
\bottomrule
\end{longtable}}
